The integration of agentic AI-systems capable of autonomous reasoning, planning, and action represents a transformative shift in cybersecurity operations ~\cite{lazer2026surveyagenticaicybersecurity}. 
However, the dual-use nature of these technologies necessitates structured engineering frameworks to ensure effective human-agent teaming. 
Evertsz and Thangarajah propose a BDI-based methodology that uses diagrammatic representations to present contextualized team cognition to humans, facilitating their role as peers rather than mere supervisors of artificial agents \cite{evertsz2020framework}. 
A primary barrier to successful teaming is the "black-box" nature of deep learning models, which often leads to alert fatigue and a lack of trust among security analysts. 
Recent research emphasizes Explainable AI (XAI) \cite{kalimuthu2025explainable} as the critical bridge to resolving this opacity. 
For example, studies utilizing SHAP and LIME demonstrate that providing feature-level justifications for IDS alerts can significantly improve forensic analysis and decision-making \cite{kalimuthu2025explainable}. 
Furthermore, evidence from user studies indicates that hybrid UI strategies combining confidence visualizations with structured rationales are most effective at calibrating user trust and improving decision accuracy. 
Ultimately, reliability in autonomous cybersecurity is increasingly defined as a socio-technical property. 
Adebayo’s HITL-XAI framework \cite{adebayo2026human} embeds human expertise at strategic junctures of the cyber kill chain, using XAI-driven uncertainty quantification to trigger human intervention only when the AI reaches its limits.
This collaborative paradigm is echoed in governmental principles for Operational Technology (OT), which stress that humans must remain responsible for functional safety through continuous oversight and fail-safe mechanisms. 
Together, these papers advocate for a shift from passive reporting to symbiotic collaboration, where AI amplifies human capacity while human experts provide the essential context for high-stakes reasoning.